\chapter{TỔNG QUAN ĐỀ TÀI}
\section{Bối cảnh đề tài}
% Nhu cầu trải nghiệm mùi hương trước khi mua; hạn chế của tester thủ công.
\section{Lý do chọn đề tài}
% Tính cần thiết của nền tảng trải nghiệm nước hoa tự động.
\section{Mục tiêu đề tài}
\subsection{Mục tiêu tổng quát}
% Xây dựng nền tảng IoT đa thương hiệu cho trải nghiệm và quản lý nước hoa.
\subsection{Mục tiêu cụ thể}
\begin{itemize}
\item Hỗ trợ khách chọn mùi hương và tạo yêu cầu trải nghiệm.
\item Quản lý thương hiệu, sản phẩm, máy và khe chứa.
\item Ghi nhận thanh toán, kết quả cấp phát và lịch sử giao dịch.
\item Theo dõi tồn kho ước tính và trạng thái máy.
\end{itemize}
\section{Đối tượng và phạm vi nghiên cứu}
\subsection{Đối tượng}
% Khách hàng, thương hiệu, quản trị viên và thiết bị kiosk.
\subsection{Phạm vi}
% Phân biệt chức năng triển khai, mô phỏng, mới thiết kế và ngoài phạm vi.
\section{Phương pháp thực hiện}
% Phân tích yêu cầu, mô hình hóa, thiết kế, phát triển, kiểm thử.
\section{Công nghệ và công cụ dự kiến}
% Điền công nghệ thực tế được sử dụng.
\section{Kết quả dự kiến}
% Liệt kê sản phẩm bàn giao theo phạm vi đồ án.
\section{Bố cục báo cáo}
% Tóm tắt nội dung các chương.
